% Options for packages loaded elsewhere
\PassOptionsToPackage{unicode}{hyperref}
\PassOptionsToPackage{hyphens}{url}
%
\documentclass[
]{article}
\usepackage{lmodern}
\usepackage{amssymb,amsmath}
\usepackage{ifxetex,ifluatex}
\ifnum 0\ifxetex 1\fi\ifluatex 1\fi=0 % if pdftex
  \usepackage[T1]{fontenc}
  \usepackage[utf8]{inputenc}
  \usepackage{textcomp} % provide euro and other symbols
\else % if luatex or xetex
  \usepackage{unicode-math}
  \defaultfontfeatures{Scale=MatchLowercase}
  \defaultfontfeatures[\rmfamily]{Ligatures=TeX,Scale=1}
\fi
% Use upquote if available, for straight quotes in verbatim environments
\IfFileExists{upquote.sty}{\usepackage{upquote}}{}
\IfFileExists{microtype.sty}{% use microtype if available
  \usepackage[]{microtype}
  \UseMicrotypeSet[protrusion]{basicmath} % disable protrusion for tt fonts
}{}
\makeatletter
\@ifundefined{KOMAClassName}{% if non-KOMA class
  \IfFileExists{parskip.sty}{%
    \usepackage{parskip}
  }{% else
    \setlength{\parindent}{0pt}
    \setlength{\parskip}{6pt plus 2pt minus 1pt}}
}{% if KOMA class
  \KOMAoptions{parskip=half}}
\makeatother
\usepackage{xcolor}
\IfFileExists{xurl.sty}{\usepackage{xurl}}{} % add URL line breaks if available
\IfFileExists{bookmark.sty}{\usepackage{bookmark}}{\usepackage{hyperref}}
\hypersetup{
  hidelinks,
  pdfcreator={LaTeX via pandoc}}
\urlstyle{same} % disable monospaced font for URLs
\usepackage{graphicx,grffile}
\makeatletter
\def\maxwidth{\ifdim\Gin@nat@width>\linewidth\linewidth\else\Gin@nat@width\fi}
\def\maxheight{\ifdim\Gin@nat@height>\textheight\textheight\else\Gin@nat@height\fi}
\makeatother
% Scale images if necessary, so that they will not overflow the page
% margins by default, and it is still possible to overwrite the defaults
% using explicit options in \includegraphics[width, height, ...]{}
\setkeys{Gin}{width=\maxwidth,height=\maxheight,keepaspectratio}
% Set default figure placement to htbp
\makeatletter
\def\fps@figure{htbp}
\makeatother
\setlength{\emergencystretch}{3em} % prevent overfull lines
\providecommand{\tightlist}{%
  \setlength{\itemsep}{0pt}\setlength{\parskip}{0pt}}
\setcounter{secnumdepth}{-\maxdimen} % remove section numbering

\date{}

\begin{document}

\hypertarget{header-n2}{%
\section{- Nov 26 Tue}\label{header-n2}}

\hypertarget{header-n3}{%
\subsection{SE-302::Compilers}\label{header-n3}}

這一部分內容就進入了「第二部分」：所謂高級主題。

今天要講的東西，我們不會實現。

就是所謂的「Garbage Collection」，垃圾回收。

\hypertarget{header-n7}{%
\subsubsection{GC: Mark \& Sweep}\label{header-n7}}

\hypertarget{header-n8}{%
\paragraph{Collect What?}\label{header-n8}}

對於 C/C++ 類語言（我們最熟悉的）來說，局部變量會被放在 Stack 上；而
\texttt{xalloc} family 開的內存會被放在 Heap 中。

局部變量所佔用的空間在編譯期基本可以確定（除了 C 的 Variable Array
這類東西），所以做垃圾回收基本是對 Heap 來做。

\hypertarget{header-n11}{%
\paragraph{Algorithms}\label{header-n11}}

\begin{quote}
MARK \& SWEEP
\end{quote}

標記 + 清掃算法。

\begin{figure}
\centering
\includegraphics{/Users/yue/Documents/GitHub/2019-2020-autumn-semester/11-november/26.assets/image-20191126081008568.png}
\caption{}
\end{figure}

一般認為：棧上的變量為根節點，也就是說認為棧所引用的對象（至少在當前）是不需要被回收的。

進行一次
Traverse，標記好所有被引用的節點，然後從內存中掃（Sweep）掉所有沒被 Mark
的對象。

\hypertarget{header-n18}{%
\paragraph{Side Effects}\label{header-n18}}

要進行 M\&S，程序需要從正常執行中中斷。

對於 Multi-Thread 程序，每個 Thread 都需要被同時打斷。

\hypertarget{header-n21}{%
\paragraph{Pseudo-Code}\label{header-n21}}

\begin{figure}
\centering
\includegraphics{/Users/yue/Documents/GitHub/2019-2020-autumn-semester/11-november/26.assets/image-20191126081459777.png}
\caption{}
\end{figure}

用一個 DFS（深度優先搜索）做 Mark 標記。

\hypertarget{header-n24}{%
\paragraph{Time Cost}\label{header-n24}}

\begin{figure}
\centering
\includegraphics{/Users/yue/Documents/GitHub/2019-2020-autumn-semester/11-november/26.assets/image-20191126081600943.png}
\caption{}
\end{figure}

活對象的數目記為 \(R\)；堆的總大小（活對象和死對象的總體）記為 \(H\)。

總時間複雜度基本在 \(\frac {c_1R + c_2H} {H - R}\)。

當 \(R\) 非常接近 \(H\) 時，也就是大部分堆上的對象都是 Alive
的時候，時間複雜度是很爆炸的。

因此這提醒了我們：沒啥垃圾的時候別亂 GC。

\hypertarget{header-n30}{%
\subsubsection{M \& S: Optimization}\label{header-n30}}

\hypertarget{header-n31}{%
\paragraph{Explicit Stack}\label{header-n31}}

在作者那個年代，調用深度一深就容易 StackOverflow⋯

因此他提供了一個優化方法：自己實現了一個 Stack Frame 的模擬，以免 DFS
過深的時候調用堆棧爆炸。

\begin{figure}
\centering
\includegraphics{/Users/yue/Documents/GitHub/2019-2020-autumn-semester/11-november/26.assets/image-20191126082156767.png}
\caption{}
\end{figure}

現在沒什麼必要了⋯⋯

\hypertarget{header-n36}{%
\paragraph{Pointer Reversal}\label{header-n36}}

指針反轉法：

\begin{figure}
\centering
\includegraphics{/Users/yue/Documents/GitHub/2019-2020-autumn-semester/11-november/26.assets/image-20191126082639844.png}
\caption{}
\end{figure}

重用堆上的數據結構，確保棧上佔用的空間不爆炸。

多線程的時候，可能就會 Conflict，或者產生一系列的 Race 之類的。

\hypertarget{header-n41}{%
\subsubsection{GC: Reference Count}\label{header-n41}}

另一種垃圾回收算法。

\hypertarget{header-n43}{%
\paragraph{Algorithms}\label{header-n43}}

使用內存中的另一塊來存放一個表，意在記錄每一個對象被引用的次數，即
Reference Count。

每次一個對象引用了 Object 的指針的時候，就將 RefCount 加
1；釋放的時候，將 RefCount 減 1。

RefCount 減為 0，代表其可以被安全刪除。

\hypertarget{header-n47}{%
\paragraph{Glitch: Cycle Referencing}\label{header-n47}}

相互引用的物體，是無法被回收的。

循環引用是 RefCount 法最大的問題。

這個問題只能由程序員解決；實現上不允許兩個對象互相引用。

\hypertarget{header-n51}{%
\subsubsection{\texorpdfstring{\textasciitilde{}Break\textasciitilde{}}{\textasciitilde Break\textasciitilde{}}}\label{header-n51}}

基本上的思路就分為兩種：

要麼用 Mark \& Sweep 找出活著的對象，然後掃掉剩餘的；

要麼用 Ref Count 找出死了的對象，單獨把他們殺掉。

有別的方法嗎？

\hypertarget{header-n56}{%
\subsubsection{Copying Collection}\label{header-n56}}

\hypertarget{header-n57}{%
\paragraph{Idea}\label{header-n57}}

\begin{figure}
\centering
\includegraphics{/Users/yue/Documents/GitHub/2019-2020-autumn-semester/11-november/26.assets/image-20191126085729115.png}
\caption{}
\end{figure}

把堆分成兩部分；一部分是 \texttt{from-space}，另一部分是
\texttt{to-space}。

我們定期將 Object 在兩塊堆之中進行拷貝；但是，僅僅通過 Traverse
的方式拷貝那些活對象。拷貝完成之後，交換 \texttt{from} 和 \texttt{to}
的地位。

\begin{quote}
類似 OpenGL 雙緩衝的技術實現。
\end{quote}

這樣的成本要低很多；他用的是 Mark \& Sweep 的思路，但是無需進行 Sweep。

而且對多線程來說也更友好；不用 Hang
整個程序；拷貝到一半的時候，也就是指針還沒被完全重建的時候也不影響正常讀取。

\hypertarget{header-n65}{%
\paragraph{Implementation}\label{header-n65}}

\begin{figure}
\centering
\includegraphics{/Users/yue/Documents/GitHub/2019-2020-autumn-semester/11-november/26.assets/image-20191126090534899.png}
\caption{}
\end{figure}

\hypertarget{header-n67}{%
\subsection{SE-227::CSE}\label{header-n67}}

\end{document}
